\begin{table}[htbp]

\centering

\resizebox{\linewidth}{!}{
\begin{tabular}{llll}
\toprule
Feature source & Entity & Feature & Description \\
\midrule
\multicolumn{4}{l}{\textbf{LP-relaxation features}} \\
\midrule
LP solver
& Variable $i$
& $c_i$
& Objective coefficient / variable weight. \\

LP solver
& Variable $i$
& $\mathbbm{1}\{x_i^{\mathrm{LP}}=0\}$
& Indicator that the LP relaxation assigns variable $i$ to zero. \\

LP solver
& Variable $i$
& $\mathbbm{1}\{x_i^{\mathrm{LP}}=1\}$
& Indicator that the LP relaxation assigns variable $i$ to one. \\

LP solver
& Variable $i$
& $\mathbbm{1}\{0 < x_i^{\mathrm{LP}} < 1\}$
& Indicator that variable $i$ is fractional in the LP relaxation. \\

LP solver
& Variable $i$
& $\operatorname{basis}_i$
& One-hot encoding of the LP simplex basis status. \\

LP solver
& Variable $i$
& $\bar{c}_i$
& Reduced cost of variable $i$ in the LP relaxation. \\

LP solver
& Variable $i$
& $x_i^{\mathrm{LP}}$
& Value of variable $i$ in the LP relaxation. \\

LP solver
& Constraint/factor $F$
& $s_F$
& Slack of the corresponding LP constraint. \\

LP solver
& Constraint/factor $F$
& $\mathbbm{1}\{s_F=0\}$
& Indicator that the constraint is active in the LP relaxation. \\

LP solver
& Constraint/factor $F$
& $\lambda_F$
& Dual value associated with the constraint. \\

\midrule
\multicolumn{4}{l}{\textbf{Problem and incidence-graph features}} \\
\midrule
% Problem data
% & Variable $i$
% & $w_i$
% & Objective coefficient, also available directly from the problem instance. \\

Problem data
& Constraint/factor $F$
& $\cos(2)$
& Constant factor-type feature. \\

Incidence graph
& Edge $(i,\FACTOR)$
& $a_{i,\FACTOR}$
& Constraint coefficient, normalized per constraint. \\

Incidence graph
& Node $v$
& $\log(1+\deg(v))$
& Log-transformed degree in the variable--constraint graph. \\

Incidence graph
& Node $v$
& $\operatorname{RWPE}_{16}(v)$
& Random-walk positional encoding with 16 dimensions. \\

Incidence graph
& Node $v$
& $\operatorname{LapPE}_{16}(v)$
& Laplacian eigenvector positional encoding with 16 dimensions. \\

Random initialization
& Node $v$
& $\xi_v \in \mathbb{R}^{10}$
& Random node features concatenated to the GNN input. \\

\bottomrule
\end{tabular}
}

\caption{Input features used by the learning models, separated by source: original problem data, LP-relaxation information, and structural features of the variable--constraint graph.}
\label{tab:dataset_features}
\end{table}